%Seitengroesse, Schriftgroesse und Vorlage (unbedingt noetig)
\documentclass[a4paper, 12pt]{article}

%Benoetigte Pakete werden (zur Uebersicht) in einer anderen Datei geladen
\input{src/packages}

%Gestaltungen (auch in anderer Datei moeglich)
\fancyhf{} %Überschreibt den Kopfzeilen- und Fusszeilenstil der Vorlage
\rfoot{\thepage} %Seitenzahl unten rechts
\renewcommand{\headrulewidth}{0pt} %Loescht die Linie in der Kopfzeile aus der Vorlage


%Hier endet der Vorspann (fuer seitenuebergreifende Entscheidungen)
\begin{document}

%Das Titelblatt wird zur Uebersicht in einer anderen Datei gestaltet
\input{src/titlepage}

%%%%%%%%%%%%%%%%%%%%%%%%%%%%%
%%%Ab hier Inhalt einfügen%%%
%%%%%%%%%%%%%%%%%%%%%%%%%%%%%


%In dem Design, das fuer eine Abschnittsueberschrift festgelegt ist, wird dies (mit automatischer Nummerierung) erzeugt
\section{Einleitung} 

%label: in Zielname, zu dem gesprungen werden kann
\subsection{Begriffsklärung} \label{sec:begriffsklaerung} 

\subsubsection*{Die Methode der Polonaise} %Abschnitte mit * erscheinen nicht im Inhaltsverzeichnis
Mit \emph{Polonaise} ist in dieser Abhandlung ausdrücklich nicht der polnische Paar-Tanz gemeint, sondern ein Gruppentanz, inspiriert durch Gottlieb Wendehals' Lied \emph{Polonaise Blankenese}, in dem er das Publikum auffordert:

    \begin{center} %Zentriert den Abschnitt dazwischen
         \color{red}„Wir ziehen los mit ganz großen Schritten,\\und Erwin fasst der Heidi von hinten an die Schulter."\color{black} (nach \cite{FutureTVGroupGmbH.2019})
    \end{center}


Durch Ausführung der Anweisungen in diesem Stimmungslied wird eine Kette von Tänzern gebildet, die sich in Schlangenlinien durch den Raum bewegt (zum Beispiel auf Karnevalsveranstaltungen, siehe auch Abb.\ref{fig:polonaise}).

%Bilder einfuegen ueber den Assistenten
\begin{figure}[htbp] 
  \centering
     \includegraphics[width=0.7\textwidth]{literatur/img/masked-ball-1176141_960_720.jpg}
  %[Text im Abbildungsverzeichnis]{Text unter dem Bild}
  \caption[Polonaise an einer Karnevals-Veranstaltung: \url{https://cdn.pixabay.com/photo/2016/02/02/21/26/masked-ball-1176141_960_720.jpg} aufgerufen am 17.11.2019]{Polonaise an einer Karnevals-Veranstaltung}
  \label{fig:polonaise} %Name des Sprungziels
\end{figure}

%Alle hintereinander folgenden Leerzeilen in Latex ergeben nur einen Zeilenumbruch. Eine Leerzeile erhält man durch \\ mehrere Leerzeilen durch \vspace{Abstand in cm}

Den Schüler*Innen sollte dieser Tanz von Geburtstagen, Karnevalsfeiern und dem Schulball bekannt sein.\\

%TODO: Überleitung

\subsubsection*{Wichtige Aspekte der Aufmerksamkeit}

Im folgenden wird die Thematik der \emph{Aufmerksamkeit} näher betrachtet.

Die Taxonomie der Aufmerksamkeit nach Sturm (erschienen in \cite{Lehrner.2011}) kann folgendermaßen zusammengefasst werden:

%Tabellen mit Assistenten
%mit p{groesse} kann eine feste Spaltengroesse festgelegt werden
\begin{tabular}{|p{0.2\linewidth}|p{0.3\linewidth}|p{0.4\linewidth}|}
	\hline 
	% inhalt spalte 1 & inhalt spalte 2 & inhalt spalte 3 \\
	Dimension & Bereich & Paradigmen \\ 
	\hline 
	Intensität & Aufmerksamkeits-aktivierung (Alertness) & einfache visuelle und auditorische Reaktionszeitaufgaben; mit und ohne Warnreiz \\ 
	\cline{2-3} %nur eine linie in spalte 2 bis 3
	 & Daueraufmerksamkeit & langandauernde Signalentdeckungsaufgaben; Signale häufig \\ 
	\cline{2-3}
	 & Vigilanz & langandauernde Signalentdeckungsaufgaben; Signale selten\\
	 \hline
	 räumliche Aufmerksamkeit (Orienting) & räumliche Verschiebung des Aufmerksamkeits-fokus & Wechsel der räumlichen Aufmerksamkeit \\
	 \hline
	 Selektivität & selektive oder fokussierte Aufmerksamkeit & Wahlreaktionsaufgaben; Reagieren auf Signale vor dem Hintergrund von Störreizen\\
	 \cline{2-3}
	  & geteilte Aufmerksamkeit & Dual-task-Aufgaben; Multi-task-Aufgaben\\
	  \hline
\end{tabular} 

%TODO: genauer erklären wie Entspannung und Lockerung dazu beitragen kann

\subsection{Beschreibung der Versuchsgruppe}

In Absprache mit Herrn Langsam habe ich die Polonaise in der Klasse 8f durchgeführt.

%TODO: Herrn Langsam zur Klasse fragen

\subsection{Begründung des Versuchs}

%TODO: Recherchieren zu Bewegung und Konzentration

%Fuer einen bewussten Seitenumbruch
\newpage

\section{Praxisteil}

\subsection{Durchführung einer Polonaise im Unterricht}

%Bilder neben Text zu positionieren funktioniert über direkt aufeinanderfolgende minipages, die nebeneinander dargestellt werden

%80 Prozent der Zeile soll aus Text bestehen

	\begin{figure}[htb]
	\begin{minipage}[t]{0.7\textwidth}
		\vspace{0pt} %Damit beide minipages oben beginnen
		Zu Beginn hat Herr Langsam die Klasse informiert, dass wir sie an diesem Tag gemeinsam unterrichten. Das Thema der Stunde war der Umfang eines Kreises (siehe Abb. \ref{fig:umfang}):\\
		
		\begin{center}
			$ U=2\cdot r\cdot \pi $ für positive Radien $ r $\\
		\end{center}
	
		\vspace{0.5cm}
		 
		 Für meinen Einstieg nutzte ich einen Filmausschnitt zur Verfilmung der Lektüre "Reise um die Erde in 80 Tagen" (siehe \cite{Verne.2015}), welche die Klasse zu dem Zeitpunkt im Französisch-Unterricht las.
	\end{minipage}
	\hspace{0.5cm}
	\begin{minipage}[t]{0.2\textwidth}
		\vspace{0pt} %Damit beide minipages oben beginnen
		\centering
		\includegraphics[width=\textwidth]{literatur/img/umfang} 
		\caption[Umfang eines Kreises: \url{https://www.pirabel.de/wp-content/uploads/2018/08/kreis-mittelpunkt-radius-300x297.png} aufgerufen am 17.11.2019]{Umfang eines Kreises}
		\label{fig:umfang}
	\end{minipage}
	\end{figure}

Anhand dessen wurde diskutiert, wie viele Schritte man braucht, um einmal um die Welt zu gehen. Nach einer Ideensammlung an der Tafel erarbeitete die Lehrkraft die Formel und übte sie an geeigneten Beispielen.\\
Nach einer Übungsphase habe ich das Polonaise-Lied abgespielt (siehe Kapitel \ref{sec:begriffsklaerung}).

%TODO: Durchfuehrung beschreiben

Zum Abschluss der Polonaise-Phase habe ich die Schüler*Innen in einen Kreis geführt, um an einem Praxisbeispiel den Umfang eines Kreises zu berechnen.

%TODO: Durchfuehrung beschreiben

\subsection{Feedback und Reflexion}

%TODO: Feedbackbögen auswerten, Herrn Langsam zu Lernerfolg fragen

Das einprägsamste Zitat war folgendes:

%Rahmen um einen Absatz durch fbox um eine minipage
\fbox{\begin{minipage}{\linewidth}
		Das war eine Mathe-Stunde, die werde ich nie vergessen!
\end{minipage}}

\section{Fazit}

%TODO

%%%%%%%%%%%%%%%%%%%%%%%%%%%%%%%%%%%%%%%%%%%%%%%%%%%
%%%%%Ende des "normalen" Inhalts des Dokuments%%%%%
%%%%%%%%%%%%%%%%%%%%%%%%%%%%%%%%%%%%%%%%%%%%%%%%%%%
\newpage
%Zur Uebersicht ausgelagert
\input{src/endofdocument}
\end{document}